\documentclass[11pt]{article}

\usepackage[margin=1in]{geometry}
\usepackage[T1]{fontenc}
\usepackage[utf8]{inputenc}
\usepackage{booktabs}
\usepackage{longtable}
\usepackage{array}
\usepackage{xcolor}
\usepackage{hyperref}
\usepackage{listings}
\usepackage{amsmath}
\usepackage{siunitx}

\hypersetup{
  hidelinks,
  pdftitle={GF/PED Manual},
  pdfauthor={Merlino Development Team}
}

\definecolor{MerlinoBlue}{HTML}{245C73}
\definecolor{MerlinoSlate}{HTML}{4A5561}
\definecolor{MerlinoGold}{HTML}{C88A2D}
\definecolor{MerlinoBg}{HTML}{F7F8F6}
\emergencystretch=2em

\lstdefinestyle{merlino}{
  basicstyle=\ttfamily\small,
  breaklines=true,
  columns=fullflexible,
  keepspaces=true,
  frame=single,
  rulecolor=\color{MerlinoSlate!45},
  backgroundcolor=\color{MerlinoBg},
  xleftmargin=0.5em,
  xrightmargin=0.5em
}

\newcommand{\program}{\texttt{GF/PED}}
\newcommand{\merlino}{\texttt{Merlino}}
\newcommand{\gicforge}{\texttt{GICForge}}
\newcommand{\bohr}{\text{bohr}}

\title{\textbf{GF/PED Manual}\\
\large Wilson GF analysis, internal-coordinate force constants and potential-energy distributions in Merlino}
\author{Merlino Development Notes}
\date{\today}

\begin{document}
\maketitle

\begin{abstract}
\program{} is the \merlino{} workflow that transforms a Cartesian harmonic
Hessian into an internal-coordinate Wilson GF problem, solves the harmonic
vibrational eigenproblem and reports frequencies, internal force constants,
normal modes and potential-energy distributions (PEDs).  The production path
uses a frozen non-redundant \gicforge{} definition, so vibrational analysis,
Gaussian input generation and structural refinement can share the same
coordinate model.  This manual describes the theory, input contracts, command
line and GUI workflows, Pulay-style scaling, output files and validation rules.
\end{abstract}

\tableofcontents
\newpage

\section{Purpose}

The GF/PED module answers a specific question: given a Cartesian Hessian and a
non-redundant internal-coordinate basis, what are the harmonic frequencies,
normal modes, internal force constants and coordinate contributions in that
basis?  It is deliberately separate from VPT2/VCI.  VPT2/VCI works in
normal-coordinate quartic force fields; GF/PED works with a Cartesian Hessian,
a Wilson B matrix and a non-redundant coordinate model.

The preferred production workflow is \texttt{gic-gf}.  It reads a frozen
\gicforge{} schema and a Gaussian FCHK Hessian adapter, evaluates the same GICs
on the Hessian geometry or on an optional external geometry, transforms the
Cartesian Hessian to the GIC basis and solves Wilson GF.  A simpler diagnostic
workflow, \texttt{gf}, reads an FCHK file and generates a Merlino GIC basis on
the fly; it is useful for quick checks but does not provide the same
cross-module coordinate identity as a frozen \gicforge{} schema.

\section{Theoretical Background}

Let \(x\) be the \(3N\)-dimensional Cartesian coordinate vector, \(F_x\) the
Cartesian harmonic Hessian in \(E_h/\bohr^2\), and \(m_i\) the atomic masses in
amu.  For a non-redundant internal coordinate vector \(q\), the linearized
relation between Cartesian and internal displacements is
\[
  dq = B\,dx ,
\]
where \(B=\partial q/\partial x\) is the Wilson B matrix.  The kinetic-energy
matrix in internal coordinates is
\[
  G = B\,M^{-1}\,B^T ,
\]
where \(M^{-1}\) is the diagonal Cartesian inverse-mass matrix with each atomic
mass repeated for \(x,y,z\).

The internal-coordinate Cartesian backtransform used by the code is
\[
  A = M^{-1} B^T G^+ ,
\]
where \(G^+\) is the numerical inverse or pseudoinverse of \(G\).  The internal
force-constant matrix is then
\[
  F_q = A^T F_x A .
\]
The matrix is explicitly symmetrized before the GF solution.  If Pulay-style
diagonal scaling factors \(s_i\) are supplied, the scaled internal Hessian is
\[
  (F_q^\mathrm{scaled})_{ij}
  =
  (F_q)_{ij}\sqrt{s_i s_j}.
\]
This preserves symmetry and applies off-diagonal scaling as the geometric mean
of the two diagonal factors.

The Wilson GF eigenproblem is solved in symmetric form.  Since \(G\) is
positive definite for a valid non-redundant basis, the code diagonalizes
\[
  G^{1/2} F_q G^{1/2}.
\]
The eigenvalues are converted to signed harmonic frequencies in cm\(^{-1}\).
Negative eigenvalues are retained as negative frequencies and therefore expose
transition-state or instability directions rather than being silently clipped.

The internal normal modes are back-expressed in the GIC basis, and the
potential-energy distribution for coordinate \(i\) in mode \(k\) is evaluated
from the mode vector \(l_k\) as
\[
  \mathrm{PED}_{ik}
  =
  100\,
  \frac{|l_{ik}(F_q l_k)_i|}
       {\sum_j |l_{jk}(F_q l_k)_j|}.
\]
The absolute-value normalization makes the PED table robust for coupled modes
with positive and negative cross terms, while retaining the relative
coordinate contribution pattern.

\section{Input Contracts}

\subsection{Cartesian Hessian}

The canonical internal object is \texttt{merlino\_gf.HessianInput}.  It stores
atomic numbers, Cartesian coordinates in bohr, atomic masses in amu, a full
symmetric Cartesian Hessian in \(E_h/\bohr^2\), optional harmonic frequencies
from the source adapter and provenance metadata.  The current file adapter is
Gaussian FCHK/FCH.

The FCHK file must contain:

\begin{itemize}
  \item atomic numbers;
  \item current Cartesian coordinates;
  \item atomic masses;
  \item Cartesian force constants;
  \item harmonic frequencies when available.
\end{itemize}

\noindent
GF/PED does not read geometries from Gaussian log text and does not use a
Z-matrix.  The FCHK adapter is only a source of canonical Cartesian Hessian
data.

\subsection{Frozen GIC Definition}

The production branch reads a \texttt{merlino.gic.definition.v1} JSON schema
created by \texttt{gic-define}.  The schema contains primitive coordinates, the
GIC coefficient matrix, labels, names, irreducible representations, point group,
symmetrization flag and the Gaussian-readable GIC block.  GF/PED evaluates that
frozen definition; it does not rerun topology perception, pruning or symmetry
assignment.

The atom count in the schema and in the FCHK input must agree.  If an optional
geometry is supplied for the B matrix, it must describe the same atom ordering
and the same topology.  Coordinates supplied through \texttt{--geometry} are in
angstrom and are converted internally to bohr before B-matrix evaluation.

\section{Command-Line Workflows}

\subsection{Quick FCHK workflow}

The quick workflow builds a Merlino GIC basis directly from the FCHK geometry:

\begin{lstlisting}[style=merlino]
python -m merlino gf \
  --fchk working/gauin.fchk \
  --out working/gf_ped_report.txt \
  --csv-dir working/gf_csv
\end{lstlisting}

\noindent
This branch is useful for diagnostics.  Because the GICs are generated inside
the GF run, it should not be used when the analysis must be identical to a
previous Gaussian input, SEfit/MORPHEUS refinement or GICForge manifest.

\subsection{Frozen-GIC production workflow}

The production workflow first freezes the coordinate model and then runs GF/PED:

\begin{lstlisting}[style=merlino]
python -m merlino gic-define \
  --geometry parent.xyz \
  --out working/gic_definition.json \
  --gaussian-out working/gauin

python -m merlino gic-gf \
  --schema working/gic_definition.json \
  --fchk working/gauin.fchk \
  --out working/gic_gf_ped_report.txt \
  --csv-dir working/gic_gf_csv
\end{lstlisting}

\noindent
If the B matrix must be evaluated on a geometry different from the FCHK
geometry, pass:

\begin{lstlisting}[style=merlino]
python -m merlino gic-gf \
  --schema working/gic_definition.json \
  --fchk working/gauin.fchk \
  --geometry working/current.xyz \
  --out working/gic_gf_ped_report.txt
\end{lstlisting}

\noindent
This is useful for controlled diagnostics, but production vibrational analysis
should normally evaluate the B matrix on the same equilibrium geometry used for
the Hessian.

\section{Pulay-Style Scaling}

GF/PED accepts optional diagonal scaling factors in the GIC basis.  Scaling is
applied after transformation of the Cartesian Hessian to internal coordinates
and before the Wilson GF eigenproblem.  A scaling file is line-oriented text or
CSV:

\begin{lstlisting}[style=merlino]
# selector factor
default 1.000
GIC003 0.980
A1Str0001,0.995
5=1.020
\end{lstlisting}

\noindent
Selectors may be:

\begin{itemize}
  \item \texttt{default}, \texttt{all} or \texttt{*};
  \item a one-based index such as \texttt{5};
  \item a \texttt{GICnnn} label;
  \item an exact GIC name;
  \item an exact label;
  \item an unambiguous substring of a name or label.
\end{itemize}

\noindent
Inline records are supplied with repeated \texttt{--scale} options:

\begin{lstlisting}[style=merlino]
python -m merlino gic-gf \
  --schema working/gic_definition.json \
  --fchk working/gauin.fchk \
  --scale-file working/pulay_scale.txt \
  --scale "GIC003=0.980" \
  --scale "default=1.000"
\end{lstlisting}

\noindent
Scaling factors must be non-negative.  Ambiguous or unmatched selectors are
fatal errors because silent scaling of the wrong coordinate would invalidate the
PED interpretation.

\section{Output Files}

\begin{longtable}{@{}p{0.31\linewidth}p{0.59\linewidth}@{}}
\toprule
Output & Meaning \\
\midrule
\endhead
\texttt{gf\_ped\_report.txt} & Human-readable report for the quick
\texttt{gf} workflow. \\
\texttt{gic\_gf\_ped\_report.txt} & Human-readable report for the frozen-GIC
production workflow. \\
\texttt{gf\_manifest.json} & Run manifest for the quick workflow. \\
\texttt{gic\_gf\_manifest.json} & Run manifest recording FCHK, schema,
optional geometry, scaling input and output files. \\
\texttt{\detokenize{*_frequencies.csv}} & Mode index and frequency in cm\(^{-1}\). \\
\texttt{\detokenize{*_gic_labels.csv}} & GIC index, name, irrep and coordinate label. \\
\texttt{\detokenize{*_ped.csv}} & Potential-energy distribution, rows as GICs and columns
as modes. \\
\texttt{\detokenize{*_normal_modes.csv}} & Internal-coordinate normal-mode vectors. \\
\texttt{\detokenize{*_force_constants.csv}} & Internal-coordinate force-constant matrix
after optional scaling. \\
\texttt{\detokenize{*_g_matrix.csv}} & Wilson G matrix in the same GIC order. \\
\bottomrule
\end{longtable}

The text report begins with the source FCHK, coordinate source, point group,
symmetrization flag and GIC count.  It then prints frequencies, GIC labels and
the PED matrix.  In the frozen-GIC workflow, labels and irreps are those stored
or evaluated from the \gicforge{} schema.

\section{GUI Workflow}

The advanced GUI exposes GF/PED through the dedicated \texttt{GF / PED} window.
The window reads a Cartesian Hessian FCHK/FCH file, an optional frozen GIC
definition, an optional B-matrix geometry and optional Pulay scaling records.
It displays the text report and can export the same CSV tables written by the
CLI.

The dashboard also exposes a \texttt{GIC Definition / B Matrix / GF-PED}
workflow.  The intended sequence is:

\begin{enumerate}
  \item define or load a frozen GIC schema;
  \item select the FCHK Hessian;
  \item optionally select a current geometry for B evaluation;
  \item optionally supply Pulay scaling;
  \item run GF/PED and export report/CSV tables.
\end{enumerate}

\section{Validation and Failure Policy}

GF/PED must fail rather than continue when the mathematical model is invalid.
Fatal conditions include:

\begin{itemize}
  \item FCHK missing masses, Cartesian coordinates or Cartesian force constants;
  \item non-symmetric Cartesian Hessian;
  \item schema atom count different from the Hessian input atom count;
  \item B matrix dimension inconsistent with the Hessian;
  \item non-positive-definite Wilson G matrix for the selected coordinates;
  \item negative Pulay scaling factor;
  \item ambiguous or unmatched Pulay selector.
\end{itemize}

\noindent
A negative vibrational eigenvalue is not a fatal error by itself.  It is
reported as a negative frequency because it may represent a transition-state
direction or an instability in the supplied Hessian/geometry combination.

\section{Relation to GICForge, MORPHEUS and VPT2/VCI}

\gicforge{} owns coordinate construction, pruning and symmetry labels.
GF/PED consumes that frozen coordinate model and performs harmonic vibrational
linear algebra.  MORPHEUS/SEfit consumes the same coordinate model for
semiexperimental structure refinement.  VPT2/VCI is a separate workflow based
on normal-coordinate cubic and quartic force fields and should not be confused
with GIC-based Wilson GF/PED.

This separation is intentional.  A coordinate definition can be generated once,
used to prepare Gaussian input, reused for semiexperimental refinement, and
then reused again for GF/PED analysis of the Cartesian Hessian.  When the same
schema is used throughout, coordinate labels, symmetry species, B matrices and
PED assignments remain auditable across modules.

\end{document}
